 \begin{document}

\title{$\MIP^* = \RE$}
\author[1]{Zhengfeng Ji\footnote{Email: zhengfeng.ji@uts.edu.au}}
\author[2,3]{Anand Natarajan\footnote{Email: anandn@caltech.edu}}
\author[3]{Thomas Vidick\footnote{Email: vidick@caltech.edu}}
\author[2,3,4]{John Wright\footnote{Email: wright@cs.utexas.edu}}
\author[5]{Henry Yuen\footnote{Email: hyuen@cs.toronto.edu}}
\affil[1]{Centre for Quantum Software and Information, University of Technology Sydney}
\affil[2]{Institute for Quantum Information and Matter, California Institute of Technology}
\affil[3]{Department of Computing and Mathematical Sciences, California Institute of Technology}
\affil[4]{Department of Computer Science, University of Texas at Austin}
\affil[5]{Department of Computer Science and Department of Mathematics, University of Toronto}

\date{\today}
\maketitle

\noteswarning


\begin{abstract}
We show that the class $\MIP^*$ of languages that can be decided by a
classical verifier interacting with multiple all-powerful
quantum provers sharing entanglement is equal to the class $\RE$ of recursively
enumerable languages. Our proof builds upon the quantum low-degree test of (Natarajan and Vidick, FOCS 2018) and the classical low-individual degree test of (Ji, et al., 2020) by integrating recent developments from (Natarajan and Wright, FOCS 2019) and combining them with the recursive compression framework of (Fitzsimons et al., STOC 2019).

An immediate byproduct of our result is that there is an efficient reduction from the Halting Problem to the problem of deciding whether a two-player nonlocal game has entangled value $1$ or at most $\frac{1}{2}$.
Using a known connection, undecidability of the entangled value implies a negative answer to Tsirelson's problem: we show, by providing an explicit example, that the closure $C_{qa}$ of the set of quantum tensor product correlations is strictly included in the set $C_{qc}$ of quantum commuting correlations. Following work of (Fritz, Rev. Math. Phys. 2012) and (Junge et al., J. Math. Phys. 2011) our results provide a refutation of Connes' embedding conjecture from the theory of von Neumann algebras.

\end{abstract}

\newpage

\tableofcontents

